\begin{remark}[E10: what PromQL does with a series that stops]\label{rem:fill}
Two rows of the beyond-reach column carried a proof and two carried a sentence, and the
two sentences covered more occurrences than the two proofs. This is the first of the two
readings that fixed that, and it did not go the way the sentence predicted.

The claim was that no PromQL term introduces a sample where the database holds none. The
setup that tests it is a gauge with two series, \texttt{svc="a"} throughout and
\texttt{svc="b"} stopping halfway, backfilled into Prometheus~3.15.0 and read back across
the whole span. Prometheus carries a stopped series forward for five minutes, so the
question is asked ten minutes past the stop, where anything present is the expression's
doing and not the staleness lookback's.

\begin{center}\footnotesize
\begin{tabular}{@{}llc@{}}
\toprule
Candidate & Expression & \texttt{svc="b"} past the lookback \\
\midrule
bare selector & \texttt{m} & absent \\
the one the argument considered & \texttt{m or vector(0)} & absent; a bare \texttt{\{\}} appears \\
matched union & \texttt{m or on() vector(0)} & absent \\
absence as a value & \texttt{m or (absent(m) * 0)} & absent \\
held alive & \texttt{last\_over\_time(m[W])} & present, last value \\
held alive and zeroed & \texttt{m or (last\_over\_time(m[W]) * 0)} & present, $0$ \\
\bottomrule
\end{tabular}
\end{center}

\noindent
So the argument was right about \texttt{or vector(0)} --- it supplies one unlabelled
series for the whole expression, which is why that candidate was rejected --- and wrong
to generalise from it. The last two rows do per-series what \texttt{.fill(last)} and
\texttt{.fill(zero)} do, with the series' own labels, which is the thing the row was
said to be impossible.

Three qualifications, none of which rescues the original claim. The fill lasts exactly
as long as $W$ covers the stopped series' last sample and no longer, so $W$ is doing the
work the dashboard's time range does in Datadog --- the row is translatable only
relative to context the query does not carry, which is \S\ref{sec:caveats}'s subject
arriving from a direction this paper did not expect. Multiplying by zero drops
\texttt{\_\_name\_\_}, so the filled series is not label-identical to the one it stands
in for; whether that matters depends on what the panel groups by. And this is a reading
of what the server does, not a soundness proof: it moves the row out of
\emph{inexpressible}, which is a claim about the language, and into \emph{unbuilt},
which is a claim about us.

The corpus makes the correction cheap to state and awkward to have needed. All $36$
occurrences are zero-fills but one --- $26$ \texttt{.fill(zero)}, $9$ \texttt{.fill(0)},
$1$ \texttt{.fill(last)} --- so there is no interpolating case among them at all, and
the row's name has been a poor fit for its contents since it was written
(\S\ref{sec:gaps} says so about a different row for the same reason). Linear
interpolation, the construct the name describes, is not in this corpus and this
experiment says nothing about it.

That makes five refusals this paper has read as a property of the languages and had to
return as a property of its own translator, and like the other four it moved the number
in the direction that made the paper's thesis sound stronger.
\end{remark}

\begin{remark}[E11: a pattern where a label name goes]\label{rem:label-name}
The other argued row, and this one survives. A Datadog scope term with no key ---
\texttt{\{*-location,*-name,*-uid,*\}}, which is $24$ of the $28$ distinct cases, and
\texttt{*-zone}, which is the other $4$ --- globs over the whole \texttt{key:value}
string rather than over a value, so expressing it in PromQL needs a pattern in a
matcher's \emph{name} position.

This is a claim about syntax, so our own grammar settles most of it: in
\S\ref{sec:grammars}'s PromQL rules a matcher's name is a terminal, never an expression,
which is why \texttt{regular.py} is never asked the question. But the grammar is
inferred, and the worry that makes a server worth asking is specific: 3.x added quoted
UTF-8 label names, and a quoted name looks syntactically like a place a pattern could
go. Five queries, sent to the same 3.15.0 server:

\begin{center}\footnotesize
\begin{tabular}{@{}lll@{}}
\toprule
Attempt & Query & Server \\
\midrule
pattern in name position & \texttt{\{=\~{}".*-zone"\}} & parse error \\
pattern as a name, with a value & \texttt{\{\~{}".*-zone"="aws"\}} & parse error \\
quoted pattern as a name & \texttt{\{".*-zone"="aws"\}} & accepted, $0$ series \\
control: quoted literal name & \texttt{\{"svc"="a"\}} & accepted, $1$ series \\
control: regex in value position & \texttt{\{svc=\~{}"a.*"\}} & accepted, $1$ series \\
\bottomrule
\end{tabular}
\end{center}

\noindent
The third row is the informative one and the reason the experiment was worth running:
the quoted form is accepted and returns nothing, because it asks for a label whose name
is literally \texttt{.*-zone}. The extension opened the name position to more
\emph{characters}, not to patterns. With the controls showing that quoted names and
value-position regexes both work on this server, what is left is that the name position
takes a literal and only a literal, which is what the row claimed.
\end{remark}
